\begin{question}

Suppose a $2 \times 2$ matrix $X$ (not necessarily Hermitian, nor unitary) is
written as:

\begin{equation*}
  X = a_0 + \vec{\sigma} \cdot \vec{a}
\end{equation*}

where the matrices $\vec{\sigma}$ are given in (3.50) and $a_0$ and
$a_{1,2,3}$ are numbers.

\begin{questionpart}
  How are $a_0$ and $a_k$(k=1,2,3) related to $tr(X)$ and
  $tr(\sigma_kX)$?
\end{questionpart}

\begin{questionpart}
  Obtain $a_0$ and $a_k$ in terms of the matrix elements $X_{ij}$
\end{questionpart}

\end{question}

\begin{answer}
Before jumping into the mathematical detail, it's worth taking a
moment to consider the physical interpretation of $\vec{\sigma} \cdot
\vec{a}$.  In this case, $\vec{a}$ is a vector representing a quantum
state where it tells you what fraction of the state is in $x$, what
fraction is in $y$, and what fraction is in $z$.  Since a has no
restriction in magnitude, presumably its magnitude is $\frac \hbar 2$
or something.  The next problem uses $\hat{n}$ as a unit vector to
much the same effect and will be discussed further in that problem.

\begin{equation*}
  \tag{3.50}
  \sigma_1 =
  \begin{pmatrix}
    0 & 1 \\
    1 & 0 \\
  \end{pmatrix},
  \sigma_2 =
  \begin{pmatrix}
    0 & -i \\
    i &  0 \\
  \end{pmatrix},
  \sigma_3 =
  \begin{pmatrix}
    1 &  0 \\
    0 & -1 \\
  \end{pmatrix}
\end{equation*}

\begin{equation*}
  \tag{3.56}
  \begin{split}
    \vec{\sigma} \cdot \vec{a} &\equiv \sum_k a_k \sigma_k \\
    &= \begin{pmatrix}
      a_3        & a_1 - ia_2 \\
      a_1 + ia_2 & -a_3 \\
    \end{pmatrix}
  \end{split}
\end{equation*}


\begin{answerpart}{}

  \begin{equation}
    \begin{split}
      tr(X)
      &= tr(a_0 \cdot I) + tr\left( \sum_i a_i \sigma_i \right) \\
      &= 2a_0 + \sum_i a_i tr(\sigma_i) \\
      &= 2a_0 + \sum_i a_i \cdot 0 \\
      &= 2a_0
      \qed
    \end{split}
  \end{equation}

  \begin{equation}
    \forall_k \  tr(\sigma_k) = 0
  \end{equation}
  
  \begin{equation}
    \inlabel{a-start}
    \begin{split}
      tr(\sigma_k X)
      &= tr\left( a_0\sigma_k + \sum_i{a_i\sigma_i\sigma_k} \right) \\
      &= a_0tr(\sigma_k) + tr\left( \sum_i{a_i\sigma_i\sigma_k} \right) \\
      &= tr\left( \sum_i{a_i\sigma_i\sigma_k} \right) \\
      &= \sum_i{tr(a_i\sigma_i\sigma_k)} \\
    \end{split}
  \end{equation}

  Let's look at the various options for $\sigma_i\sigma_j$.  For
  $i=j$, we recall equation equation 3.51a:

  \begin{equation*}
    \tag{3.51a}
    \sigma_i^2 = I
  \end{equation*}

  \begin{equation}
    tr(\sigma_i\sigma_i) = tr(I) = 2
  \end{equation}

  \begin{equation*}
    \tag{1.172a}
    tr(XY) = tr(YX)
  \end{equation*}

  There are, in principle, 9 combinations.  3 of those are for $i=j$.
  Of the remaining 6, 3 are equivalent to the other 3 via equation
  1.172a (e.g. $\sigma_x\sigma_y = \sigma_y\sigma_x$).  Of the
  remaining 3, two of them involve $\sigma_z$ ($\sigma_z\sigma_x$ and
  $\sigma_z\sigma_y$) both of which are trivially 0 owing to the fact
  that $\sigma_z$ is diagonal and $\sigma_x/\sigma_y$ are both
  anti-diagonal.

  The final case is $\sigma_x\sigma_y$:
  \begin{equation}
    \begin{split}
      tr(\sigma_x\sigma_y)
      &=
      tr\left(
      \begin{bmatrix}
        0 & 1\\
        1 & 0\\
      \end{bmatrix}
      \begin{bmatrix}
        0 & -i\\
        i & 0\\
      \end{bmatrix} \right)\\
      &=
      tr\left(\begin{bmatrix}
        i & 0\\
        0 & -i\\
      \end{bmatrix} \right)\\
      &= 0
    \end{split}
  \end{equation}

  That allows us to generalize the results with:
  
  \begin{equation}
    tr(\sigma_i\sigma_j) = 2\delta_{ij}
  \end{equation}

  Substituting our new identity into equation \inref{a-start}
  
  \begin{equation}
    \begin{split}
      tr(\sigma_k X)
      &= \sum_i{tr(a_i\sigma_i\sigma_k)} \\
      &= \sum_i{2a_i\delta_{ik}} \\
      &= 2a_k
      \qed
    \end{split}
  \end{equation}

\end{answerpart}


%%%%%%%%%%%%%%%%%%%%%%%%%%%%%%%%%%%%%%%%%%%%%%%%%%%%%%%%%%%%%%%%%%%%%%%%%%%%%%%
%%                              Part b                                       %%
%%%%%%%%%%%%%%%%%%%%%%%%%%%%%%%%%%%%%%%%%%%%%%%%%%%%%%%%%%%%%%%%%%%%%%%%%%%%%%%

\begin{answerpart}{}

  This is mostly already provided in the book in eqaution 3.56

  \begin{equation*}
    \tag{3.56}
    \begin{split}
      \vec{\sigma} \cdot \vec{a} &\equiv \sum_k a_k \sigma_k \\
      &= \begin{pmatrix}
        a_3        & a_1 - ia_2 \\
        a_1 + ia_2 & -a_3 \\
      \end{pmatrix}
    \end{split}
  \end{equation*}

  Adding in $a_0$
  
  \begin{equation}
    \inlabel{solution-b}
    \begin{split}
      \vec{\sigma} \cdot \vec{a} &\equiv \sum_k a_k \sigma_k \\
      &= \begin{pmatrix}
        a_0 + a_3  & a_1 - ia_2 \\
        a_1 + ia_2 & a_0 - a_3 \\
      \end{pmatrix}
      \qed
    \end{split}
  \end{equation}

\end{answerpart}
\end{answer}
